% Parked 2026-08-13 (author directive: attribution section commented out).
% Restore by splicing back into results.tex before sec:dense_weak.

\subsection{Per-Bucket Attribution}
\label{sec:buckets}

The reference point for the attribution is the per-bar least-squares
fit of the HAR ladder plus calendar dummies, QLIKE \unifAzeroQLIKE{}.
The first measurement is two quantities: the change in QLIKE when each
exogenous feature bucket is added to that backbone, and the overlap
between the buckets' individual contributions. The buckets are defined
in Section~\ref{sec:dp_data}. The estimator is held fixed at the
least-squares class of Section~\ref{sec:estimators} throughout the
attribution; the same features under richer hypothesis classes are
treated in Sections~\ref{sec:dense_weak} and \ref{sec:tree_edge}.

Table~\ref{tab:buckets} reports the causally tuned ridge arm on each of the
exogenous buckets, on the shared panel, against the shared incumbent.

% Provenance: writeup/metrics_table_causal_tune_plus_spectral.csv
% (family=causal-tuned, estimator ridge), incumbent from
% writeup/incumbent_metrics.json; Δ/Δ_own columns as in
% reference/sections/marginal_contribution.tex Table tab:buckets.
\begin{table}[H]
\centering
\caption{Bucket attribution: minimum-norm least squares on each exogenous
bucket, from the unification campaign. \textbf{Panel:} the frozen panel of
Section~\ref{sec:data_prep}, $n = 273{,}554$ thirty-minute bars, identical
rows in every cell. \textbf{Protocol:} rows are scored under the campaign's
the evaluation protocol of
Section~\ref{sec:smearing_contract}; the incumbent every $\Delta$ and
DM in this table is computed against is QLIKE \unifAzeroQLIKE{}. (The
first-iteration version of this table, scored under the scalar protocol,
is superseded; the text records what changed and why.)
\textbf{Estimator:} minimum-norm least squares (Section~\ref{sec:estimators}),
rolling window $24{,}000$ bars ($500$ trading days), refit every bar. Columns
participate in a window only when their feed went live before the window
starts and printed within it; masked columns are disclosed per chunk.
$\Delta$ is against the incumbent; the within-class increment coincides
with $\Delta$ because the no-exogenous design is exactly the incumbent.}
\label{tab:buckets}
\begin{tabular}{lrrr}
\toprule
Bucket & QLIKE & $\Delta$ vs.\ incumbent & DM $t$ \\
\midrule
all\_features (joint) & \unifAllFeaturesQLIKE{} & \unifAllFeaturesDelta{} & \unifAllFeaturesDM{} \\
\midrule
moments      & \unifMomentsQLIKE{}    & \unifMomentsDelta{}    & \unifMomentsDM{} \\
liquidity    & \unifLiquidityQLIKE{}  & \unifLiquidityDelta{}  & \unifLiquidityDM{} \\
implied\_vol & \unifImpliedVolQLIKE{} & \unifImpliedVolDelta{} & \unifImpliedVolDM{} \\
vol\_demand  & \unifVolDemandQLIKE{}  & \unifVolDemandDelta{}  & \unifVolDemandDM{} \\
sentiment    & \unifSentimentQLIKE{}  & \unifSentimentDelta{}  & \unifSentimentDM{} \\
market\_ew   & \unifMarketEwQLIKE{}   & \unifMarketEwDelta{}   & \unifMarketEwDM{} \\
market\_vw   & \unifMarketVwQLIKE{}   & \unifMarketVwDelta{}   & \unifMarketVwDM{} \\
\midrule
(no exogenous $=$ incumbent) & \unifAzeroQLIKE{} & --- & --- \\
\bottomrule
\end{tabular}
\end{table}

On the full 1998--2024 panel, no single-bucket arm improves on the
incumbent. Two are statistically indistinguishable from it --- moments
(\unifMomentsDelta{}, $t = $ \unifMomentsDM{}) and sentiment
(\unifSentimentDelta{}, \unifSentimentDM{}) --- and the remaining five
are significantly worse: market EW (\unifMarketEwDM{}), market VW
(\unifMarketVwDM{}), implied volatility (\unifImpliedVolDM{}),
volatility demand (\unifVolDemandDM{}), and liquidity
(\unifLiquidityDelta{}, \unifLiquidityDM{}). The joint arm --- all
exogenous columns in one minimum-norm fit --- is significantly worse
still (\unifAllFeaturesDelta{}, \unifAllFeaturesDM{}), worse than every
single bucket it contains: even the strongest unpenalized estimator
extracts nothing from the stacked design, and pays the variance cost of
its width. The information these columns carry under shrinkage
(Sections~\ref{sec:dense_weak_twoblock}--\ref{sec:three_block_results})
is invisible to the unshrunk fit. That contrast --- information present,
unpenalized extraction null, shrunk extraction strong --- is the table's
central finding, and the first-iteration table it replaces told the same
story more weakly: under the earlier scalar protocol the two best
buckets appeared to win outright, an artifact the current model-specific
second moment removes --- the correction is largest exactly where the
mean model is weakest, which is the benchmark, not the bucket arms.

A disclosure attaches to the market-aggregate rows. In earlier passes of
this campaign both aggregates appeared catastrophically harmful. The cause
was a single evaluation window (December 2016 -- February 2017) in which a
pair of availability-indicator columns, differing only through pre-2015
history, decayed to exact collinearity as that history left the rolling
window; the linear solver does not signal a nearly singular system, and the
contaminated coefficients produced forecast errors that one window's
correction then spread across all of its bars. The per-bar solve
verification of the final design (Section~\ref{sec:data_prep}) closes this
failure mode; under it the aggregates score as reported in
Table~\ref{tab:buckets}.

\subsection{Bucket Attribution under Penalized Estimators}
\label{sec:buckets_tuned}


% Provenance: 16 arms (br_tuned_*, be_tuned_*), each a full 100-chunk run on
% the frozen panel; causally tuned ridge (logspace(-2,3,6)) and causally
% tuned elastic net with FREE mixing (alpha logspace(-6,-2,5) x l1_ratio
% {0.25,0.5,0.75,1.0}), TUNE_PER=250, 25-bar embargo, 125-bar tail.
% Selected (alpha, l1_ratio) persisted per retune. 2026-08-07.
Table~\ref{tab:buckets} answers what an unpenalized estimator recovers.
It is not the same question as what the columns contain. Each bucket
design is therefore re-run under two penalized estimators, both choosing
their penalty by the causal rule of Section~\ref{sec:estimators}: ridge,
and an elastic net whose mixing parameter is \emph{free} --- the tuner
selects $\ell_1$ share as well as penalty strength at every retune, so
the amount of selection is chosen by the estimator rather than fixed by
the author. Table~\ref{tab:buckets_tuned} reports both.

\begin{table}[H]
\centering
\small
\caption{Each bucket design under two causally tuned penalized
estimators. \textbf{Panel and protocol:} as in Table~\ref{tab:buckets}
--- the frozen panel of Section~\ref{sec:data_prep}, identical rows in
every cell, rows scored under the evaluation protocol of
Section~\ref{sec:smearing_contract} (benchmark \unifAzeroQLIKE{}). DM $t$ is against the benchmark; negative
favours the arm. The final column is the paired Diebold--Mariano
statistic of the elastic net against ridge \emph{on the same design},
positive meaning ridge is the better of the two.}
\label{tab:buckets_tuned}
\begin{tabular}{@{}lrrrrr@{}}
\toprule
& \multicolumn{2}{c}{Ridge} & \multicolumn{2}{c}{Elastic net} & Enet vs.\ ridge \\
\cmidrule(lr){2-3}\cmidrule(lr){4-5}\cmidrule(l){6-6}
Design & QLIKE & DM $t$ & QLIKE & DM $t$ & DM $t$ \\
\midrule
moments      & \unifBrTunedMomentsQLIKE{}    & \unifBrTunedMomentsDM{}    & \unifBeTunedMomentsQLIKE{}    & \unifBeTunedMomentsDM{}    & \unifIncrBeTunedMomentsDM{} \\
liquidity    & \unifBrTunedLiquidityQLIKE{}  & \unifBrTunedLiquidityDM{}  & \unifBeTunedLiquidityQLIKE{}  & \unifBeTunedLiquidityDM{}  & \unifIncrBeTunedLiquidityDM{} \\
market (EW)  & \unifBrTunedMarketEwQLIKE{}   & \unifBrTunedMarketEwDM{}   & \unifBeTunedMarketEwQLIKE{}   & \unifBeTunedMarketEwDM{}   & \unifIncrBeTunedMarketEwDM{} \\
market (VW)  & \unifBrTunedMarketVwQLIKE{}   & \unifBrTunedMarketVwDM{}   & \unifBeTunedMarketVwQLIKE{}   & \unifBeTunedMarketVwDM{}   & \unifIncrBeTunedMarketVwDM{} \\
sentiment    & \unifBrTunedSentimentQLIKE{}  & \unifBrTunedSentimentDM{}  & \unifBeTunedSentimentQLIKE{}  & \unifBeTunedSentimentDM{}  & \unifIncrBeTunedSentimentDM{} \\
implied vol. & \unifBrTunedImpliedVolQLIKE{} & \unifBrTunedImpliedVolDM{} & \unifBeTunedImpliedVolQLIKE{} & \unifBeTunedImpliedVolDM{} & \unifIncrBeTunedImpliedVolDM{} \\
vol.\ demand & \unifBrTunedVolDemandQLIKE{}  & \unifBrTunedVolDemandDM{}  & \unifBeTunedVolDemandQLIKE{}  & \unifBeTunedVolDemandDM{}  & \unifIncrBeTunedVolDemandDM{} \\
\midrule
all exogenous (joint) & \unifBrTunedAllFeaturesQLIKE{} & \unifBrTunedAllFeaturesDM{} & \unifBeTunedAllFeaturesQLIKE{} & \unifBeTunedAllFeaturesDM{} & \unifIncrBeTunedAllFeaturesDM{} \\
\bottomrule
\end{tabular}
\end{table}

Two readings follow, and they motivate presenting the
attribution and the estimator comparison jointly.

First, \emph{shrinkage reverses the joint design from liability to
contender}. Under minimum-norm least squares the joint arm was
significantly worse than every bucket it contains
(Table~\ref{tab:buckets}); under a causally tuned ridge it is the
second-best row in the table (\unifBrTunedAllFeaturesQLIKE{}, DM
\unifBrTunedAllFeaturesDM{}) and, with the moments bucket
(\unifBrTunedMomentsQLIKE{}, DM \unifBrTunedMomentsDM{}), one of only
two rows not significantly worse than the benchmark. The strong claim of
the first-iteration table --- the joint design made outright best ---
does not survive the current protocol; what survives is the reversal
itself: with an estimator that shrinks, more exogenous columns stop
costing and start paying, and the single best row in the table beats the
benchmark outright. What Table~\ref{tab:buckets} measures is the cost of
forgoing shrinkage, and that cost is larger than everything the buckets
individually contain.

Second, ridge attains a lower pooled loss than the elastic net on every
information set: all eight paired comparisons favour ridge except
liquidity, which is a tie, and four do so at conventional significance
(vol.\ demand \unifIncrBeTunedVolDemandDM{}, moments
\unifIncrBeTunedMomentsDM{}, sentiment \unifIncrBeTunedSentimentDM{},
market VW \unifIncrBeTunedMarketVwDM{}). The elastic net's failure mode
is visible in the table itself: the volatility-demand arm scores
\unifBeTunedVolDemandQLIKE{} --- worse than its own unpenalized fit in
Table~\ref{tab:buckets} --- a catastrophic penalty selection of exactly
the kind the analysis below locates. That ordering, however, is not yet
a statement about the two \emph{families}, for the reason given next.

\paragraph{Where the elastic net's deficit comes from.}
% Provenance: per-bar losses of be_tuned_* and br_tuned_* partitioned by
% the (alpha, l1_ratio) the tuner selected at the governing retune;
% 9,600 retunes over 8 designs; pooled differentials reproduce
% results/unification_increments.csv to 10 decimal places. 2026-08-08.
\emph{The decomposition in this paragraph is computed on the
first-iteration protocol's arms; the mechanism it documents is what the
current protocol's volatility-demand arm exhibits live in
Table~\ref{tab:buckets_tuned}.}
Because the elastic net selects its penalty causally at every retune,
each scored bar can be attributed to the penalty in force when it was
forecast, and the deficit can be located rather than merely totalled.
It does not spread across the sample. On the $72.1\%$ of bars whose
governing penalty was interior to the grid ($\alpha \le 10^{-3}$), the
elastic net \emph{beats} tuned ridge by
$3.9 \times 10^{-4}$ (DM $-6.70$, favouring the elastic net in all eight
designs). The entire pooled deficit is contributed by the $27.9\%$ of
bars whose governing penalty sat at the \emph{top point of the grid},
$\alpha = 10^{-2}$, where it loses by $2.5 \times 10^{-3}$ (DM $+6.85$,
in all eight designs and in $18$ of $24$ years). The
difference between the two subsets is $2.9 \times 10^{-3}$ at
$z = 7.8$.

A result carried by a grid endpoint suggests that the grid may be
too small. Under the correspondence between the elastic net's
parameterization and a ridge penalty --- an $\ell_2$ weight of $N \alpha
(1 - \rho)$ on the window gram, with $N = 24{,}000$ and $\rho$ the
mixing share --- the elastic net's largest expressible $\ell_2$ weight
is $24{,}000 \times 10^{-2} \times 0.75 = 180$, while the tuned ridge
arm draws from a grid reaching $1{,}000$ and selects that top value in
$41.3\%$ of its own retunes. The elastic net cannot express the
shrinkage ridge chooses, and it loses on the bars where it has run out
of grid.

That explanation is testable, and it is rejected. Re-running every design
with the penalty grid extended to $\alpha = 1$ --- which reaches $6$ to
$18$ times ridge's selected weight at every mixing value, and contains
the original grid as a subset --- makes the elastic net
\emph{substantially worse}, not better. On the moments design the
extended-grid arm loses $2.9 \times 10^{-3}$ to its own narrow-grid twin
(DM $+7.01$) and $3.3 \times 10^{-3}$ to tuned ridge (DM $+6.68$); on
the joint design, $1.7 \times 10^{-3}$ (DM $+4.56$) and $2.4 \times
10^{-3}$ (DM $+3.46$). Giving the estimator more room to shrink causes
it to shrink worse.

This locates the asymmetry precisely, and it is not about how much
shrinkage each family can express. Both tuners pin the tops of their
grids --- ridge in $41.3\%$ of retunes, the elastic net in a comparable
share --- because on a $125$-bar validation tail neither can resolve the
penalty. What differs is the consequence. For ridge, the loss is
monotone and flat in $\log \alpha$ near its optimum
(Section~\ref{sec:dense_weak_headtohead}), so selecting too large a
penalty costs almost nothing: over-shrinking a dense coefficient vector
degrades gracefully. For the elastic net, a larger penalty is partly an
$\ell_1$ penalty, and over-penalizing in that direction removes columns
outright. The same selection error is benign in one family and
destructive in the other. Selection loses here not because the data
prefer dense coefficients, but because the price of a hyperparameter
mistake is asymmetric --- and a short causal validation window
guarantees such mistakes. This mechanism complements the asymptotic
analysis of \citet{ShenXiu2025}: their ridge-over-lasso ordering is
established under oracle or cross-validated tuning, whereas here the
ordering persists --- and is amplified --- under causal selection on a
short forward tail, whose mistakes load on exactly the dial the two
families price differently.

\paragraph{What the trajectories do and do not show.}
\emph{(Also computed on the first-iteration arms.)} The selected mixing
share is bimodal --- $36.7\%$ of retunes at pure
lasso, $32.5\%$ at the minimum $0.25$, mean $0.64$ --- and the same
shape appears in every bucket. One might read this as the estimator
declining to select; two measurements argue against that
interpretation. Holding $\alpha$ fixed, the mixing share accounts for only
$4.2 \times 10^{-4}$ of the deficit, significant in one of eight designs
and reversing sign at $\alpha = 10^{-4}$; and the selected mixing share
is unrelated to the volatility regime (Spearman $-0.013$ against the
preceding window's realized volatility, $p = 0.23$). The mixing
trajectory is a symptom of a validation tail too short to identify the
penalty, not the channel through which the loss is incurred.

Nor is ridge exempt from that instability: it pins an endpoint of its
own grid in $58.1\%$ of retunes, against $43.3\%$ for the elastic net.
The asymmetry this section can support is therefore narrower than a
family verdict --- both estimators are choosing badly on a $125$-bar
tail, and ridge's bad choices happen to be inexpensive while the
elastic net's are not.

\subsection{Overlap}
\label{sec:overlap}

The overlap ratio --- the sum of single-bucket increments divided by the
joint increment --- cannot be computed from Table~\ref{tab:buckets},
because the joint arm's increment is positive: at this dimensionality the
unshrunk joint fit's variance cost exceeds any information gain, so the
joint OLS arm does not measure the buckets' combined information. The
available overlap evidence comes from shrunk estimators: an earlier
battery under a causally tuned ridge measured $R = 2.08$ over its seven
buckets, and that battery's causally tuned LightGBM arm $R = 2.24$ ---
individual gains summing to roughly twice the joint gain, under both
estimator classes, on the prior convention. \pending{overlap ratio
recomputed from shrunk joint fits under the paper's frozen convention ---
the block estimators of Sections~\ref{sec:dense_weak_twoblock}--\ref{sec:three_block}
provide the joint fits}

\subsection{The Joint Bucket}
\label{sec:joint_bucket}

The joint \texttt{all\_features} arm --- every exogenous column in one
minimum-norm fit --- has QLIKE \unifAllFeaturesQLIKE{}, significantly
\emph{worse} than the incumbent (DM \unifAllFeaturesDM{}). The same
columns under the per-block shrinkage of Sections~\ref{sec:dense_weak_twoblock}
and \ref{sec:three_block} produce the paper's best forecasts
(\unifBlkThreeUserQLIKE{} at \unifBlkThreeUserDM{}), so what the joint
row measures is the cost of forgoing shrinkage, not the absence of
information.

The estimator choice matters here, and its history is part of the
record. Under the earlier plain normal-equations solve, this arm appeared
catastrophically worse than the incumbent; the deficit traced to
solution components in near-null directions of the design --- large
mutually canceling coefficients on nearly duplicate availability-ladder
columns, which a one-bar feed desynchronization on a forward-filled
post-early-close bar (December 24, 2020, the only such
disagreement in over twenty-five thousand bars) turned into a forecast
of $77.8$ against a target of $0.26$, whose residual the window-level
correction then spread across its whole window. The minimum-norm
solution suppresses near-null components by construction, and the
dead-session grid rule removed this bar class; under the two
repairs the arm scores as reported.

An independent every-bar robustness design was run on a different panel:
rolling $1{,}000$-trading-day window; HAR block exactly unpenalized via
Frisch--Waugh--Lovell, with only the exogenous block penalized; scored on a
fixed COVID-inclusive window from January 2018 to May 2025, $n = 74{,}934$
bars. On that panel, every bucket rejects equal predictive accuracy against
HAR-only (the weakest, \texttt{vol\_demand}, at $t = -6.2$), and the joint
set rejects equal predictive accuracy against every single bucket, at
$t = -13.0$. The bucket ranking differs across the two panels:
\texttt{implied\_vol} has the largest single-bucket increment on the
COVID-inclusive window, \texttt{moments} in Table~\ref{tab:buckets}.

An earlier draft asserted that the $90\%$ model confidence set
\citep{HansenLundeNason2011} over the feature sets is the singleton
containing the joint bucket. The per-bar loss series required to compute a
model confidence set were not retained for the battery, so no set was
computed, and the claim is not repeated here.
\pending{battery MCS over the feature sets once per-bar losses are persisted}
Table~\ref{tab:buckets} supports three statements: three buckets
individually beat the incumbent at the $5\%$ level and two more lean the
same way without significance; the unshrunk joint fit is worse than the
incumbent and than every single bucket; and extracting the buckets'
combined information therefore requires the shrinkage machinery the
remainder of the paper develops.


% Methods-side attribution protocol subsection, parked with the above.

\subsection{The Attribution Protocol}
\label{sec:attribution}

\paragraph{Protocol.} The attribution is computed as follows.
\begin{enumerate}
\item \emph{Arms.} The incumbent arm's design matrix contains the HAR ladder
columns and the calendar dummies. Each bucket arm's design matrix contains
those backbone columns plus the bucket's own columns, expanded through the
same moving-average ladder and availability indicators as the joint arm. The
joint arm's design matrix contains the backbone columns plus all $41$
exogenous series ($\approx 526$ columns after the moving-average ladder and
availability indicators).
\item \emph{Fit.} Each arm is fit by the rolling estimator of
Section~
ef{sec:estimators} and produces
one out-of-sample forecast per bar.
\item \emph{Loss.} Each arm is scored by per-bar QLIKE on the shared panel
($n = 218{,}934$ thirty-minute bars, identical rows in every cell). The
reported QLIKE is the mean of the arm's per-bar losses.
\item \emph{Increments.} $\Delta$ is the arm's QLIKE minus the incumbent's
QLIKE. $\Delta_{\text{own}}$ is the arm's QLIKE minus the same estimator's
own no-exogenous baseline ($0.13445$).
\item \emph{DM test.} For each arm, the per-bar loss differential against the
incumbent is $d_t = \ell^{\text{arm}}_t - \ell^{\text{inc}}_t$. The
Diebold--Mariano statistic \citep{DieboldMariano1995} is the $t$-ratio of
$\bar{d}$ to its Newey--West HAC standard error, with the
Harvey--Leybourne--Newbold small-sample correction
\citep{HarveyLeybourneNewbold1997}. Negative $t$ favours the arm.
\item \emph{Overlap.} The overlap ratio is
\[
R \;=\; \frac{\sum_{b} \Delta_{\text{own}}(b)}{\Delta_{\text{own}}(\text{joint})},
\]
the sum of the single-bucket within-class increments divided by the joint
arm's within-class increment, both measured against the same estimator's
no-exogenous baseline.
\end{enumerate}

